
This paper proposes a confidence geometry framework for learned termination detection in neural theorem proving. The central hypothesis is that LLM confidence variance encodes information about proof space navigation that can distinguish likely-successful from likely-divergent searches. The research presents a detailed methodology including confidence extraction, symbolic signal integration, and detector architectures.

However, critical examination reveals that the reported empirical results derive from synthetically generated mock data rather than actual theorem proving experiments. The results files contain entries explicitly labeled as mock data, and the variance values match predetermined statistical distributions specified in mock data generation code. While the research documents indicate awareness of this issue and claim it was resolved, the examined artifacts show that mock data remains in the results files.

The paper reports strong correlations (r=0.80, p<10⁻²³), large effect sizes (Cohen's d=2.21-2.65), and high detector performance (F1=0.806-0.97), but these values reflect properties of the synthetic data generation process rather than empirical findings. The paper does not clearly disclose this critical limitation to readers.

The conceptual framework may have merit and the proposed methodology appears reasonable in design. However, the empirical validation claimed in the abstract, introduction, and results sections has not been completed. The actual scientific contribution is a proposed research direction with detailed methodology, not a validated empirical finding.

Future work would need to:
\begin{enumerate}
\item Execute the experimental protocol with actual LeanDojo theorem proving
\item Collect real confidence trajectories from actual proof searches
\item Generate ground truth labels through extended-timeout experiments with real theorems
\item Measure actual detector performance on real data
\item Verify computational overhead through profiling of actual implementation
\item Complete Phase 5 baseline comparisons if empirical validation succeeds
\end{enumerate}

Until these steps are completed, the confidence geometry principle remains an untested hypothesis rather than an established finding.
